% Compact, code-aligned main-paper description of LPB DAPRO.
% A complete proof of the CRC statement is in standard_crc_budget_validity.tex.

\section{Proposed Method}
\label{sec:method}

Let $\calI=\{1,\ldots,N\}$ be the calibration indices and let $\bar B$ be
the target budget per calibration row, so the total budget is
$B_{\mathrm{tot}}:=N\bar B$.  We randomly partition $\calI$ into a fully
observed policy-fit set $\calIone$, an independent fully observed budget-control
set $\calIcrc$, and a deployment set $\calItwo$.  Write
\[
    N_1:=|\calIone|,\qquad n:=|\calIcrc|,\qquad
    m:=|\calItwo|,\qquad N=N_1+n+m.
\]
For row $i$, $X_i$ is its initial prompt, $T_i\in\{1,\ldots,\maxl+1\}$ is
its time to event, and $\maxl+1$ denotes no event by the executable horizon
$\maxl$.  Let
\begin{equation}
    Q_i:=\qprior(X_i)=\fhat_{\tauprior}(X_i),
    \qquad L_i:=\min\{T_i,Q_i\}.
    \label{eq:dapro_main_active_length}
\end{equation}
Here $\tauprior$ is the fixed prior candidate level, $Q_i$ is the largest
interaction needed by the LPB calibration candidates, and $L_i$ is the number
of interactions needed either to observe the event or to reach $Q_i$.  We
acquire all $L_i$ interactions for $i\in\calIone\cup\calIcrc$; only
$\calIone$ is used to learn the policy shape, whereas $\calIcrc$ is reserved
for budget calibration.

At turn $t$, let $H_i(t-1)$ denote the history available immediately before
deciding whether to acquire interaction $t$, and let $\Dtrain$ be the data
used to fit the survival model, independently of the calibration outcomes.
The predictable allocation score is the current event hazard
\begin{equation}
    S_i(t):=h_i(t)
    :=\widehat{\mathbb P}\!\left(
      T_i=t\mid T_i\geq t,H_i(t-1),\Dtrain
    \right).
    \label{eq:dapro_main_hazard}
\end{equation}
A policy returns a conditional continuation probability $P_i(t)\in[0,1]$.
We write
\begin{equation}
    \rho_i(t):=\prod_{s=1}^{t}P_i(s),\qquad \rho_i(0):=1,
    \qquad
    c_i(P):=\sum_{t=1}^{L_i}\rho_i(t),
    \label{eq:dapro_main_reach_and_cost}
\end{equation}
where $\rho_i(t)$ is the probability of reaching interaction $t$, conditional
on the complete trajectory, and $c_i(P)$ is the corresponding conditional
expected number of acquired interactions.  Thus $c_i(P)$ is an expected
cost---not a hard upper bound on the realized stopping time.

\paragraph{Variance-aligned target.}
For the LPB target miscoverage $\alpha$, let
$F_i^\star:=\fhat_\alpha(X_i)$ and
$A_i^\star:=\mathbb I\{T_i<F_i^\star\}$; in the implementation,
$\fhat_\alpha$ denotes the last candidate grid point strictly below $\alpha$.
If $\pi_i$ is the probability that the acquisition policy resolves the
target event for row $i$, then the fixed-candidate Horvitz--Thompson estimator
has conditional acquisition variance
\begin{equation}
    \operatorname{Var}(\widehat\mu_\alpha\mid\mathcal D)
    =\frac{1}{N^2}\sum_{i=1}^{N}
      A_i^\star\left(\frac{1}{\pi_i}-1\right),
    \label{eq:dapro_main_ht_variance}
\end{equation}
where $\mathcal D$ denotes the complete trajectories and fixed policy.  The
final calibrated level is data dependent and unknown before acquisition, so
we optimize this fixed raw-$\alpha$ target.  It is exact for that candidate
and a variance-aligned proxy for the final LPB.

Hard endpoint indicators are noisy when $N_1$ is small.  We therefore use the
causal soft-prefix target masses
\begin{equation}
    a_i(t):=
      \mathbb I\{t\leq L_i\}
      \mathbb I\{t<F_i^\star\}h_i(t),
      \qquad i\in\calIone,
    \label{eq:dapro_main_soft_mass}
\end{equation}
which replace the realized endpoint indicator by its model-estimated
conditional mass at every observed prefix.

After fully observing $\calIone$, define the remaining budget and the
policy-shape budget by
\begin{equation}
    B_{\mathrm{rem}}
      :=B_{\mathrm{tot}}-\sum_{i\in\calIone}L_i,
    \qquad
    \bar B_{\mathrm{shape}}
      :=\frac{B_{\mathrm{rem}}}{n+m}.
    \label{eq:dapro_main_shape_budget}
\end{equation}
The quantity $\bar B_{\mathrm{shape}}$ is only the fitting target used to
construct a useful policy family.  It is not the asserted deployment budget.
The full-observation cost of $\calIcrc$ is not yet known without exposing the
independent CRC outcomes and is accounted for exactly once by the CRC loss
below.

Let $M_t^{\mathrm{raw}}$ denote the learned monotone lookup from the score at
turn $t$ to a conditional continuation probability.  Soft-prefix Generalized
DAPRO fits this lookup using only $\calIone$ by solving
\begin{equation}
\begin{aligned}
    \min_{M^{\mathrm{raw}}}\quad&
      \frac{1}{N_1}\sum_{i\in\calIone}
      \sum_{t=1}^{L_i}a_i(t)
      \left\{\frac{1}{\rho_i^{\mathrm{raw}}(t)}-1\right\}\\
    \mathrm{s.t.}\quad&
      \frac{1}{N_1}\sum_{i\in\calIone}
      \sum_{t=1}^{L_i}\rho_i^{\mathrm{raw}}(t)
      \leq\bar B_{\mathrm{shape}},\\
    &S_i(t)\leq S_j(t)
      \Longrightarrow
      M_t^{\mathrm{raw}}(S_i(t))
      \leq M_t^{\mathrm{raw}}(S_j(t)),
\end{aligned}
\label{eq:dapro_main_optimization}
\end{equation}
where
$\rho_i^{\mathrm{raw}}(t):=\prod_{s=1}^{t}
M_s^{\mathrm{raw}}(S_i(s))$.  The implementation optimizes the monotone
lookup directly; it does not regress row-wise oracle probabilities onto the
scores.  Model misspecification can reduce efficiency but does not invalidate
the final HT correction, because the deployed policy is predictable and its
actual probabilities are logged.

\paragraph{Nested CRC family.}
CRC first freezes the target, score cutpoints, raw lookup, and its policy-fit
budget correction.  Let $\rho_i^{\mathrm{base}}(t)$ denote the resulting
base cumulative reach.  To obtain a finite support bound, a non-increasing
envelope $e(1),\ldots,e(\maxl)$ is fitted using only $\calIone$ and satisfies
\begin{equation}
    1\geq e(1)\geq\cdots\geq e(\maxl)\geq\epsilon,
    \qquad
    \sum_{t=1}^{\maxl}e(t)\leq C_{\max},
    \qquad
    C_{\max}:=\min\{\maxl,2\bar B\},
    \label{eq:dapro_main_pathwise_cap}
\end{equation}
where $\epsilon>0$ is a fixed minimum cumulative reach.  Define the capped
base reach
\begin{equation}
    \bar\rho_i(t):=\min\{\rho_i^{\mathrm{base}}(t),e(t)\}.
    \label{eq:dapro_main_capped_reach}
\end{equation}
Because the envelope is shared, frozen, and causal,
$\sum_{t=1}^{L_i}\bar\rho_i(t)\leq C_{\max}$ for every possible trajectory.
This bounds conditional expected cost; a realized acquisition can still
continue beyond $C_{\max}$ and as far as $Q_i$.  The cap is not
efficiency-neutral: because cumulative reach is non-increasing, for any
$q\leq Q_i$,
\begin{equation}
    \bar\rho_i(q)
      \leq\frac{1}{q}\sum_{t=1}^{q}\bar\rho_i(t)
      \leq\frac{C_{\max}}{q}.
    \label{eq:dapro_main_cap_terminal_effect}
\end{equation}
It may therefore reduce inclusion probabilities for targets resolved only at
late endpoints.  The canonical choice $C_{\max}=2\bar B$ is a design tradeoff,
not a theorem-optimal constant: increasing it enlarges the policy class but
also enlarges the finite-sample CRC correction.

For $\lambda\in[0,1]$, CRC contracts the cumulative reach, not the
conditional probabilities, according to
\begin{equation}
    \rho_i^\lambda(t)
      :=\epsilon+\lambda\{\bar\rho_i(t)-\epsilon\}.
    \label{eq:dapro_main_lambda_cumulative}
\end{equation}
The exact effect on the sequential continuation probabilities is therefore
\begin{equation}
    P_i^\lambda(1)
      =\epsilon+\lambda\{\bar\rho_i(1)-\epsilon\},
    \qquad
    P_i^\lambda(t)
      =\frac{\epsilon+\lambda\{\bar\rho_i(t)-\epsilon\}}
             {\epsilon+\lambda\{\bar\rho_i(t-1)-\epsilon\}},
      \quad t\geq2.
    \label{eq:dapro_main_lambda_conditionals}
\end{equation}
Hence the calibrated conditional policy is not an affine mixture of the raw
conditional probabilities.  Equivalently, one may write the final policy as
$M_t(H_i(t-1);\lambda):=P_i^\lambda(t)$: the raw lookup maps the current
score to the base path, while CRC calibrates that path through the single
scalar $\lambda$.  The history argument is necessary because the denominator
in~\eqref{eq:dapro_main_lambda_conditionals} contains the preceding cumulative
reach.  At $\lambda=1$ the policy equals the capped base policy;
at $\lambda=0$ it reaches every active interaction with cumulative
probability $\epsilon$---equivalently, it samples the row with probability
$\epsilon$ and then follows a sampled row through its active horizon.

For $i\in\calIcrc$, define
\begin{equation}
    b_i:=L_i,\qquad
    c_i(\lambda):=\sum_{t=1}^{L_i}\rho_i^\lambda(t),
    \qquad r:=\frac{n}{m}.
    \label{eq:dapro_main_crc_costs}
\end{equation}
Here $b_i$ is the realized cost of fully observing CRC row $i$, whereas
$c_i(\lambda)$ is the conditional expected cost that candidate $\lambda$
would have incurred on that same complete trajectory.  The composite loss
and its deterministic bound are
\begin{equation}
    \ell_i(\lambda):=c_i(\lambda)+r b_i,
    \qquad
    L_{\max}:=C_{\max}+r\maxl,
    \label{eq:dapro_main_crc_loss}
\end{equation}
so $0\leq\ell_i(\lambda)\leq L_{\max}$.

The mathematical continuous-family rule selects the most aggressive feasible
contraction,
\begin{equation}
    \widehat\lambda
    :=\sup\left\{\lambda\in[0,1]:
      \frac{
        \sum_{i\in\calIcrc}\ell_i(\lambda)+L_{\max}
      }{n+1}
      \leq\frac{B_{\mathrm{rem}}}{m}
    \right\}.
    \label{eq:dapro_main_crc_selector}
\end{equation}
The implementation evaluates the descending grid
$\Lambda_K=\{1,1-(K-1)^{-1},\ldots,0\}$, where $K$ is the number of
candidate scales, and takes its largest feasible
element, which is a conservative discretization of
\eqref{eq:dapro_main_crc_selector}; the same CRC proof applies with the
supremum replaced by the largest feasible grid element.  Since
\begin{equation}
    c_i(\lambda)
      =\epsilon L_i
       +\lambda\left\{
          \sum_{t=1}^{L_i}\bar\rho_i(t)-\epsilon L_i
       \right\},
    \label{eq:dapro_main_crc_linear_cost}
\end{equation}
the continuous value also has an explicit form.  Let
\begin{equation}
    \bar c_i:=\sum_{t=1}^{L_i}\bar\rho_i(t),\qquad
    D_{\mathrm{crc}}:=\sum_{i\in\calIcrc}
       \{\bar c_i-\epsilon L_i\},
    \label{eq:dapro_main_crc_slope}
\end{equation}
and
\begin{equation}
    R_{\mathrm{crc}}
    :=\frac{n+1}{m}B_{\mathrm{rem}}-L_{\max}
      -r\sum_{i\in\calIcrc}b_i
      -\epsilon\sum_{i\in\calIcrc}L_i.
    \label{eq:dapro_main_crc_residual}
\end{equation}
If $D_{\mathrm{crc}}>0$ and $R_{\mathrm{crc}}\geq0$, then
\begin{equation}
    \widehat\lambda_{\mathrm{cont}}
      =\min\left\{1,\frac{R_{\mathrm{crc}}}{D_{\mathrm{crc}}}\right\}.
    \label{eq:dapro_main_crc_closed_form}
\end{equation}
If $D_{\mathrm{crc}}=0$, all candidates have the same control cost and we
take $\widehat\lambda_{\mathrm{cont}}=1$ when feasible; if
$R_{\mathrm{crc}}<0$, even $\lambda=0$ is infeasible and the requested
budget and positivity floor are incompatible.

Conditional on $\calIone$, suppose that the CRC and deployment trajectories
are exchangeable, the entire nested family is frozen before observing
$\calIcrc$, $\ell_i(\lambda)$ is non-decreasing and right-continuous in
$\lambda$, and the selector is feasible.  Standard CRC then yields
\begin{equation}
    m\,\mathbb E[c_{\mathrm{new}}(\widehat\lambda)\mid\calIone]
    +n\,\mathbb E[b_{\mathrm{new}}\mid\calIone]
    \leq B_{\mathrm{rem}}.
    \label{eq:dapro_main_budget_guarantee}
\end{equation}
Here $c_{\mathrm{new}}$ and $b_{\mathrm{new}}$ are the corresponding costs
on an independent trajectory from the same conditional distribution as the
CRC and deployment rows.
Adding the already incurred policy-fit cost proves marginal expected control
of the configured total budget.  The term $r b_i$ is therefore not double
counting: it is what incorporates the unavoidable full-observation cost of
$\calIcrc$ in the same exchangeable loss.  A proof is given in
Appendix~\ref{sec:standard_crc_budget_validity}.  The result is an upper-budget
guarantee; it neither guarantees exact budget utilization nor controls every
realized split.

\paragraph{Deployment.}
For each $i\in\calItwo$, at every reached prefix we evaluate $S_i(t)$, apply
the frozen policy in~\eqref{eq:dapro_main_lambda_conditionals}, and draw the
corresponding Bernoulli continuation decision.  If $C_i$ denotes the last
interaction acquired for row $i$, the exact propensity of resolving its
active endpoint is
\begin{equation}
    \pi_i=\prod_{t=1}^{L_i}P_i^{\widehat\lambda}(t).
    \label{eq:dapro_main_logged_propensity}
\end{equation}
Rows in $\calIone\cup\calIcrc$ are fully observed and have $\pi_i=1$.
The weighted calibration curve is
\begin{equation}
    \widehat\mu(\tau)
    =\frac{1}{N}\sum_{i=1}^{N}
      \frac{
        \mathbb I\{T_i<\fhat_\tau(X_i),
                    \fhat_\tau(X_i)\leq C_i\}
      }{\pi_i},
    \label{eq:dapro_main_ht_curve}
\end{equation}
for every candidate level $\tau$, and is passed to the same strict-prefix
selector as the full-data LPB construction.
